\chapter{Arquitectura de software}

La arquitectura de la aplicación describe sus componentes internos y sus
responsabilidades. La aplicación accede a los datos a través de un único servicio,
\texttt{DataService}; no selecciona repositorios ni conoce tablas.

\section{Responsabilidades del acceso a datos}

\begin{description}
    \item[\texttt{DataService}.] Expone operaciones completas a la aplicación,
    comprueba las reglas de aplicación, coordina los repositorios implicados y
    decide cuándo confirmar (\textit{commit}) o revertir (\textit{rollback}).
    Traduce los errores de persistencia a resultados o excepciones comprensibles
    para sus consumidores.
    \item[\texttt{Repository} y sus repositorios concretos.] La interfaz común
    fija el criterio compartido. Cada repositorio conoce su tabla, columnas,
    claves y correspondencia con los modelos; realiza las consultas de su
    responsabilidad y devuelve modelos o resultados definidos. No confirma
    transacciones ni llama a otros repositorios.
    \item[\texttt{Postgres}.] Gestiona la conexión y ejecuta las operaciones
    contra PostgreSQL. Proporciona una única conexión transaccional compartida
    por los repositorios participantes y realiza el \textit{commit} o el
    \textit{rollback} que solicita \texttt{DataService}. No contiene reglas de
    aplicación ni conoce las entidades concretas.
\end{description}

Cada operación de \texttt{DataService} que modifica una o varias tablas utiliza
una sola conexión a PostgreSQL, compartida por todos los repositorios que
participan en ella. Así, sus consultas pertenecen a la misma transacción:
\texttt{DataService} solicita el \textit{commit} si todas terminan correctamente
y el \textit{rollback} si alguna falla. Se evita confirmar una parte de la
operación y dejar el resto sin aplicar. Las claves foráneas y demás restricciones
de integridad siguen siendo responsabilidad de PostgreSQL.

Los repositorios construyen las consultas específicas de sus tablas. Cuando una
operación sea común y genérica, \texttt{Postgres} puede ofrecer una primitiva
reutilizable para ejecutarla. Cada repositorio solicita directamente a
\texttt{Postgres} que ejecute sus consultas; \texttt{DataService} no transporta
SQL ni resultados del controlador.

Los repositorios no se llaman entre sí. Si una escritura referencia otra fila,
no se consulta previamente su existencia solo para repetir la garantía de la
clave foránea: PostgreSQL rechaza la escritura si la referencia no existe. El
repositorio traduce las violaciones conocidas a errores comprensibles y
\texttt{DataService} revierte la transacción y comunica el fallo. Las reglas de
negocio que no cubren las restricciones referenciales, como exigir que una fuente
esté habilitada, sí se comprueban desde \texttt{DataService} mediante los
repositorios necesarios.

La composición configura las dependencias y credenciales al iniciar la
aplicación; ese detalle no se expone a la aplicación ni a las reglas de negocio.
